\chapter{Implementació: De l'Orquestrador a les Quatre Famílies}
\label{chap:implementacio}

Aquest capítol tradueix la metodologia a programari. Primer l'orquestrador que fa possible $n=10.000$. Després cada paradigma tal com està al repositori ---no tal com el README de capçalera el va etiquetar a vegades malament. PPO es descriu com a framework, sense resultats.

\section{Mapa de software}

\begin{itemize}
    \item \texttt{src/p00\_core/engine/}: \texttt{benchmark.py} (master), \texttt{worker.py} (subprocess).
    \item \texttt{src/p00\_core/core/}: \texttt{BaseHeuristic1v1}, \texttt{AgentFactory}, matemàtica compartida.
    \item \texttt{src/p01\_heuristics/agents/internal/}: \texttt{v1.py} \ldots \texttt{v14.py}.
    \item \texttt{src/p03\_minmax/}: \texttt{v15}--\texttt{v17}.
    \item \texttt{src/p04\_mcts/}: \texttt{v18}--\texttt{v20}, \texttt{LocalSim} via fork Pokechamp.
    \item \texttt{src/p02\_imitation\_learning/}: dataset $\rightarrow$ features $\rightarrow$ XGBoost $\rightarrow$ \texttt{v21}/\texttt{v22}.
    \item \texttt{src/p05\_ppo\_drl/}: Gymnasium + MaskablePPO (\textbf{no avaluat} en aquesta memòria).
    \item \texttt{src/p00\_core/online\_bot/}: hook a la ladder pública.
\end{itemize}

No hi ha mòdul de Dobles a la memòria. El que en quedava al pla es va eliminar com a fora d'abast.

% ============================================================
\section{Arquitectura Master--Worker}
\label{sec:master_worker}

\subsection{Per què subprocessos}

Simular milers de batalles en un sol procés Python acumula objectes \texttt{Battle}, arbres de cerca i fils de \texttt{asyncio}. El GIL impedeix paral·lelisme CPU real amb threads. \texttt{fork()} copia l'event loop i mort en deadlock. La solució és \texttt{multiprocessing} amb context \texttt{spawn}: cada worker és un intèrpret nou, un port Showdown exclusiu, un CSV temporal, i en sortir el SO reclama la RAM.

\begin{figure}[H]
    \centering
    \begin{tikzpicture}[
        node distance=1cm, font=\footnotesize,
        master/.style={rectangle, draw, fill=blue!10, text width=3.8cm, align=center, minimum height=1.2cm, rounded corners},
        worker/.style={rectangle, draw, fill=green!8, text width=3cm, align=center, minimum height=1cm},
        server/.style={rectangle, draw, fill=orange!10, text width=2.8cm, align=center, minimum height=0.9cm},
        disk/.style={cylinder, draw, fill=gray!15, shape border rotate=90, aspect=0.25, minimum height=0.9cm, minimum width=2cm, align=center, font=\scriptsize}
    ]
        \node[master] (master) {\textbf{MASTER}\\\texttt{benchmark.py}};
        \node[worker, below left=2cm and 1.5cm of master] (w1) {Worker 1};
        \node[worker, below=2cm of master] (w2) {Worker 2};
        \node[worker, below right=2cm and 1.5cm of master] (wn) {Worker N};
        \node[server, below=1cm of w1] (s1) {Showdown :8000};
        \node[server, below=1cm of w2] (s2) {Showdown :8001};
        \node[server, below=1cm of wn] (sn) {Showdown :800N};
        \node[disk, right=2.5cm of wn] (tmp) {\_tmp\_*.csv};
        \node[disk, below=0.8cm of tmp] (final) {agent\_vs\_opp.csv};
        \draw[->, thick] (master) -- (w1);
        \draw[->, thick] (master) -- (w2);
        \draw[->, thick] (master) -- (wn);
        \draw[<->, thick] (w1) -- (s1);
        \draw[<->, thick] (w2) -- (s2);
        \draw[<->, thick] (wn) -- (sn);
        \draw[->, dashed, gray] (wn.east) -- (tmp.west);
        \draw[->, thick, dotted] (master.east) -| (final.north);
    \end{tikzpicture}
    \caption{Master--worker del gauntlet. Sense memòria compartida entre workers.}
    \label{fig:master_worker_full}
\end{figure}

\subsection{Resume i chunks}

Si \texttt{v12\_vs\_abyssal.csv} ja té 4.320 files, el master calcula $n_{\mathrm{to\_run}}=10.000-4.320$ i continua. No hi ha flag \texttt{--resume}: la veritat és el CSV. Cada worker juga chunks de 25, escriu, crida \texttt{reset\_battles()} i \texttt{gc.collect()}. Sense això, 10k batalles en un procés superen 2~GB; amb chunks, $\sim$200--400~MB.

Timeouts: el cost varia tres ordres (heurística $\sim$0,01--0,1~s/partida; MCTS $\sim$2--4~s). El timeout de batch és $\max(7200, 40\cdot n_{\mathrm{batch}})$ segons, per no matar un worker d'MCTS viu ni esperar eternament un penjament.

\subsection{Reinici de Showdown}

V8 fragmenta l'\textit{old generation}. Empíricament, després de milers de batalles el procés Node creix de $\sim$80~MB a $>400$~MB i la latència puja. El master fa \texttt{pkill} + relançament cada $N$ matchups actius (\texttt{--restart-every}, defecte 3). Cost $\sim$17~s de downtime, negligible en una nit.

\subsection{Esquema CSV}

Cada fila és una batalla. Identitat (\texttt{battle\_id}, \texttt{format}, \texttt{heuristic}, \texttt{opponent}, \texttt{won}, \texttt{turns}), qualitat de decisió (fallbacks, errors), HP i fainted, canvis voluntaris/forçats, RNG (crit, miss, SE), i telemetria de paradigma (\texttt{search\_diff}, \texttt{xgb\_stays}, \texttt{ko\_guards}, \texttt{terastallized}, \ldots). \texttt{StatsBattle} intercepta \texttt{|-crit|} i companyia abans que poke-env les descarti. Els agents no cal que sàpiguen que existeix.

La Figura~\ref{fig:ram_usage} (de la fase de Singles) mostra el perfil de dents de serra: pujada durant un matchup, caiguda al reinici.

\begin{figure}[H]
    \centering
    \includegraphics[width=0.85\linewidth]{ram_usage_singles1.png}
    \caption{RAM durant un benchmark de Singles. Les dents de serra són reinicis de Showdown + GC, no un memory leak «resolt» a mà dins el worker.}
    \label{fig:ram_usage}
\end{figure}

\subsection{Contracte \texttt{BaseHeuristic1v1}}

Patró Template Method: \texttt{choose\_move} crida \texttt{\_pre\_move\_hook} (sortida primerenca: KO), \texttt{\_select\_action} (lògica) i, si cal, \texttt{choose\_random\_move} (fallback). Una excepció no gela la batalla: es compta \texttt{error\_moves} i es juga atzar. Tots els agents interns hereten d'aquí, inclosos minimax, MCTS i IL.

\texttt{AgentFactory.create("v20")} és l'únic punt d'instanciació. Baselines: wrappers locals de \texttt{random}, \texttt{max\_power}, \texttt{simple\_heuristic} (còpia amb Tera, perquè el SimpleHeuristics de poke-env no n'havia), \texttt{safe\_one\_step} (reimplementació sense LocalSim per evitar penjades), i \texttt{abyssal} del fork Pokechamp.

% ============================================================
\section{Heurístiques \texttt{v1}--\texttt{v14}: ablació de coneixement}
\label{sec:heur_impl}

L'escala no és «catorze bots cada cop un xic millors». És un experiment: quina peça de coneixement mou el WR a Random Battle. La genealogia de \emph{codi} (Figura conceptual) té branques: \texttt{v1}$\to$\texttt{v2}$\to$\texttt{v3}$\to$\texttt{v6}; \texttt{v4}/\texttt{v5} a part; \texttt{v7} reescriu; \texttt{v8}--\texttt{v14} creixen d'aquest nucli estratègic. \texttt{v14} hereta \texttt{v13}, que hereta l'esperit de \texttt{v12}, no necessàriament el WR.

\paragraph{\texttt{v1}--\texttt{v6} (altiplà de dany).}
\texttt{v1} maximitza $\mathrm{bp}\times\mathrm{eff}\times\mathrm{STAB}$ i gairebé no canvia. \texttt{v2} afegeix stats i cremada. \texttt{v3} tracking de moviments i pivot (Toxic, outspeed). \texttt{v4}/\texttt{v5} clima, terrain, boosts. \texttt{v6} és un \texttt{v3} lleuger amb prioritat. El resultat empíric ---tot el clúster a 44--46\%--- és la tesi d'aquesta fase: més fidelitat a la fórmula de dany no guanya Random Battles.

\paragraph{\texttt{v7} (salt posicional).}
Reescriptura: hazards, setup, pre-check de KO, switch per puntuació de matchup (fórmula d'estil Abyssal). És el primer agent que \emph{tria} canviar, no només desmaia.

\paragraph{\texttt{v8} / \texttt{v10}.}
Meta: items (Life Orb, Choice), habilitats (Levitate, Flash Fire), pantalles, Trick Room. Status (Toxic, Will-O-Wisp, Thunder Wave), sacrifici $\le 20\%$ HP (no gastar un torn per salvar un moribund), U-turn / Volt Switch. Increments reals i petits.

\paragraph{\texttt{v9} / \texttt{v11} (tempo).}
La regla que val més que tota la taula d'objectes de \texttt{v8}: hazards i setup \textbf{només en torn lliure} (outspeed i resist STAB rival). \texttt{v11} híbriditza \texttt{v9}+\texttt{v10} i a penes mou l'agulla.

\paragraph{\texttt{v12} (Gen~9 jugada com a Gen~9).}
Tres peces natives del format:

\begin{enumerate}
    \item Teracristal·lització ofensiva i defensiva.
    \item Lead de \texttt{teampreview}: millor matchup mitjà contra les sis espècies conegudes.
    \item Switch-in després de faint per matchup, no l'slot 0 (ni \texttt{choose\_random\_move}, que és el que fan Abyssal/SimpleHeuristic quan \texttt{active is None}).
\end{enumerate}

A més conserva torns lliures, sack, pivots i KO. Empíricament és el sostre bot-contra-bot.

\paragraph{\texttt{v13}.}
Lazy-load de \texttt{sets.json}: predicció de moviments, item i ability. Matchup amb dany simulat, no només tipus. Choice-lock, phaze vs sweepers, recuperació $<60\%$ HP, Tera conservador (no gastar-lo a $<30\%$ HP ni en un status). Corregeix un fallback de switch-in de \texttt{v12}.

\paragraph{\texttt{v14} (orientat a humans).}
Yomi~2: perfil \texttt{PREDICTIVE} vs \texttt{CONSERVATIVE} segons la taxa de canvis del rival. Scouting als torns 1--3 (Protect, U-turn, Knock Off). 16 rolls de dany per KO garantit. Solver 1-ply d'endgame quan queden $\le 2$ mons per bàndol. Tera defensiu de parany. Contra un bot determinista, scouting i Yomi són sovint \emph{over-respect}: gasten tempo que \texttt{v12} gasta atacant. Això és disseny, no un bug, i el gauntlet ho quantifica.

Tots registren comptadors. \texttt{ko\_checks} explota a $\sim$15/partida a \texttt{v14} perquè el calculador de 16 rolls dispara constantment.

% ============================================================
\section{Minimax 1-ply \texttt{v15}--\texttt{v17}}
\label{sec:mm_impl}

Els tres hereten \texttt{HeuristicV14}. Abans de la matriu: KO garantit (i, a \texttt{v16}/\texttt{v17}, stop de setup, absorb de status, U-turn primerenc). Si no hi ha drecera, s'enumeren les accions pròpies contra les respostes \emph{predites} del rival (moviments revelats, omplerts amb la DB de sets, més un switch hipotètic si el matchup és dolent). Velocitat i prioritat es resolen analíticament: si el ràpid KO, l'acció del lent és nul·la. \textbf{No hi ha \texttt{LocalSim}.}

Fulla de \texttt{v15}: HP i matchup després del torn, amb dany rival $\times 1{,}5$. \texttt{v16} afegeix bonuses de setup / hazard / recuperació / status ($0{,}25$--$0{,}35$) ---la hipòtesi de l'efecte horitzó «s'arregla a la fórmula». \texttt{v17} és \texttt{v16} més $+0{,}15$ al score worst-case de l'acció que hauria triat \texttt{v14}.

Per què no 2-ply: $9\times 9=81$ cel·les a 1-ply, $81^2=6561$ a 2-ply, inassumible amb 16 rolls en Python dins el timer de Showdown. El README del repositori que anomena \texttt{v16} «2-ply» és \textbf{incorrecte}; \texttt{v16} és 1-ply amb fulla posicional.

Telemetria: \texttt{search\_fire\_rate}, \texttt{search\_diff} (override de la sonda \texttt{v14}). Sense aquests camps es podria vendre l'agent com a «minimax cada torn». Amb ells, es veu que la matriu corre en el 16--18\% dels torns.

% ============================================================
\section{IS-MCTS \texttt{v18}--\texttt{v20}}
\label{sec:mcts_impl}

Mateix professor \texttt{v14}, mateix pressupost:

\begin{itemize}
    \item 100 simulacions per torn de cerca, $C=1{,}4$.
    \item Rollout de 5 torns amb \texttt{LocalSim} (Python in-process, fork Pokechamp; no el Node.js per cada pas).
    \item Cada simulació determiniza el rival amb la DB de sets.
    \item Fills arrel: moviments \emph{i} canvis; es juga l'argmax de visites.
\end{itemize}

\texttt{v18}: UCB1, fulla HP/boosts/status/hazards, només KO abans de l'arbre. \texttt{v19}: UCB1, fulla amb rols i amenaça OHKO, i les dreceres tàctiques de \texttt{v14} \emph{abans} de l'arbre. \texttt{v20}: \texttt{v19} + \textbf{PUCT} amb prior $0{,}70$ sobre l'acció \texttt{v14}.

\texttt{LocalSim} $\neq$ Showdown. Un desajust de mecànica és soroll al backup d'UCB1 i una causa plausible de l'override-i-perdre de \texttt{v18}/\texttt{v19}. El minimax 1-ply \textbf{no} usa LocalSim: el seu override-i-perdre no es pot atribuir al simulador, sinó al maximin.

Etiqueta: \texttt{v18}--\texttt{v20} són MCTS, no IL. Documents vells del repositori que diuen el contrari es descarten.

% ============================================================
\section{Imitació \texttt{v21}--\texttt{v22}}
\label{sec:il_impl}

\subsection{Dades}

Replays humans \texttt{gen9randombattle}, Elo $\ge 1800$, HuggingFace / Pokechamp, finestra 2023--2025. $\sim$1,12 milions de torns. \texttt{GroupShuffleSplit} per \texttt{battle\_id}: els torns d'una mateixa batalla no filtren al test. Desbalanceig move/switch $\approx 2{,}88$; \texttt{scale\_pos\_weight} ho compensa.

No és OU. El model antic entrenat en OU i avaluat en Random Battle ($\sim$2\% WR) era un mismatch de format, no una refutació de l'IL.

\subsection{Cap macro}

XGBoost, 1.150 features (HP, hazards, Tera flags, one-hot d'espècie activa, \ldots), $n_{\mathrm{estimators}}=200$, $\mathrm{depth}=6$, $\eta=0{,}1$. Sortida: $p(\mathrm{switch})$. Llindar calibrat $\tau=0{,}5525$ (\texttt{xgboost\_advanced\_threshold.json}). Accuracy de classificació $\sim$65,5\%, ROC-AUC $\sim$0,72: n'hi ha prou per un prior de timing, no per un jugador.

\subsection{Dues piles}

\paragraph{\texttt{v21} (\texttt{HeuristicV21XGBoost}).}
Hereta \texttt{v14}. Ordre: KO garantit $\to$ endgame $\le 2$ mons $\to$ setup/absorb $\to$ XGBoost. Si SWITCH: pivot o \texttt{\_get\_best\_switch} de \texttt{v14}. Si MOVE: \texttt{\_score\_move} de \texttt{v14}. Guard de loop. Empíricament XGB dispara en el 15\% dels torns.

\paragraph{\texttt{v22} (\texttt{HeuristicV22PureIL}).}
Només \texttt{BaseHeuristic1v1}. XGB primer. Switch: puntua cada mons del banc com si fos l'actiu i tria el de $p(\mathrm{switch})$ mínima. Move: segon model \texttt{xgboost\_move\_evaluator.json} (atributs: power, STAB, eff, status, prioritat --- no el nom del moviment). Zero KO guards. Loop guard gairebé sempre.

El README que intercanvia híbrid i pur és \textbf{incorrecte}. \texttt{v21} és l'híbrid; \texttt{v22} és el clon.

% ============================================================
\section{Baselines}

\texttt{random}, \texttt{max\_power}: poke-env. \texttt{simple\_heuristic}: còpia local amb Tera. \texttt{abyssal}: Pokechamp. \texttt{one\_step} i \texttt{safe\_one\_step} apunten al player matemàtic sense LocalSim. Abyssal i Simple Heuristic són, en aquest gauntlet, gairebé el mateix rival (Elo 1755 tots dos, empatats amb \texttt{v14}).

% ============================================================
\section{Bot en línia}

\texttt{run\_online\_bot.py} connecta qualsevol etiqueta de l'\texttt{AgentFactory} al servidor públic. Mode \texttt{ladder} o \texttt{accept}. La mostra de tesi és \texttt{v14}, usuari \texttt{SirPThesis}, 9--10 de juny de 2026, 431 partides. No s'ha fet un ladder de \texttt{v12}: el protocol bot-contra-bot i el protocol humà es queden separats.

% ============================================================
\section{PPO (framework, no experiment d'aquesta memòria)}
\label{sec:ppo_impl}

\texttt{src/p05\_ppo\_drl/} encapsula un \texttt{PokemonMaskedEnv} (accions 0--3 moviments, 4--9 switches), vectorització, MaskablePPO~\cite{sb3_contrib} i un currículum Random $\to$ MaxPower $\to$ SimpleHeuristic $\to$ gauntlet. Això documenta que el cinquè paradigma \emph{es pot} entrenar sobre el mateix Showdown. \textbf{No s'hi reporta cap WR.} La comparació prevista és vs \texttt{v12}, a part, no una fila més de la matriu 28.
